\documentclass[11pt,a4paper]{article}
\usepackage[utf8]{inputenc}
\usepackage[margin=25mm]{geometry}
\usepackage[T1]{fontenc}
\ifdefined\DriftJaxCAS\else
\usepackage{lmodern}
\fi
\usepackage{amsmath,amssymb,mathtools}
\usepackage{graphicx}
\graphicspath{{records/}}
\usepackage{booktabs,tabularx,array}
\usepackage{xcolor}
\usepackage[expansion=false]{microtype}
\usepackage{listings,enumitem}
\usepackage[numbers,sort&compress]{natbib}
\usepackage{hyperref}
\hypersetup{colorlinks=true,linkcolor=blue!45!black,
  citecolor=blue!45!black,urlcolor=blue!45!black}
\lstset{basicstyle=\ttfamily\small,breaklines=true,language=Python,
  frame=single,backgroundcolor=\color{gray!4},showstringspaces=false,
  columns=fullflexible,keepspaces=true}
\setlist{itemsep=2pt,topsep=4pt}
\setlength{\emergencystretch}{2em}
\newcommand{\code}[1]{\texttt{\detokenize{#1}}}
\newcommand{\dd}{\mathrm{d}}
\newcommand{\jac}{\mathcal{J}}
\newcommand{\design}{\boldsymbol{\theta}}
\newcommand{\state}{\boldsymbol{u}}
\newcommand{\res}{\boldsymbol{F}}
\newcommand{\papertitle}{DriftJax: A Structure-Aware Differentiable
Drift--Diffusion Solver for Photovoltaic Device Design}
\newcommand{\paperauthors}{Mohammed Benaissa$^{a}$\textsuperscript{*} \and
Mohamed Maoudj$^{a}$ \and Abd Elouahab Noua$^{a,b}$}
\newcommand{\paperaffiliations}{%
\begin{center}\small
$^{a}$Research Center in Semiconductor Technology for Energetics (CRTSE),
02 Bd. Frantz Fanon, B.P. 140, 16038 Algiers, Algeria\\
$^{b}$Department of Sciences and Technology, University Larbi Ben M'hidi,
04000 Oum El Bouaghi, Algeria\\[3pt]
\textsuperscript{*}Corresponding author: Mohammed Benaissa.
\end{center}}
\newcommand{\paperfigure}[3]{%
\ifdefined\DriftJaxCAS
\begin{figure}[pos=tbp]
\else
\begin{figure}[tbp]
\fi
\centering
\includegraphics[width=0.96\linewidth]{#1}
\caption{#2}\label{#3}
\end{figure}}
\newenvironment{papertable}{%
\ifdefined\DriftJaxCAS
\begin{table}[pos=tbp]
\else
\begin{table}[tbp]
\fi
}{\end{table}}
\usepackage{caption}
\captionsetup{font=small,labelfont=bf}
\renewcommand{\thesection}{S\arabic{section}}
\numberwithin{equation}{section}
\renewcommand{\thefigure}{S\arabic{figure}}
\renewcommand{\thetable}{S\arabic{table}}
\title{\large Supplementary material\\[6pt]\Large\papertitle}
\author{\paperauthors}
\date{}
\begin{document}
\maketitle
\paperaffiliations

\noindent This supplement accompanies both the standard and Elsevier CAS
versions of the article. Sections S1--S6 develop the mathematical
conventions and verification procedures; S7 presents additional device
examples; S8 records the provenance of retained numerical values; and S9
specifies the information required for reproduction. Derivations are
distinguished from implementation claims and from reported measurements.
Implementation descriptions follow the v0.1.17 source tree described in
the main article; per-example JSON sidecars, execution logs, the
benchmark record (\code{bench\_v017\_dgbsv.json}), and the heterojunction
gradient-evidence record (\code{gradient\_evidence\_hetero.json}) are
archived under \code{docs/paper/records/}. Statements scoped to the
v0.1.16 revision below are retained as historical provenance.

\section{Physical conventions and nondimensionalisation}\label{supp:physics}
\subsection{State variables and units}
The device occupies $0\leq x\leq d$. Electron and hole electrical current
densities are positive along $+x$, and the elementary charge $q$ is
positive. In these conventions, the standard transport equations
are~\cite{gaury2019,mann2022partial}
\begin{align}
\partial_x(\varepsilon\partial_x\phi)
 &=q(n-p+N_A^- -N_D^+),\\
\partial_xJ_n&=q(R-G),&\partial_xJ_p&=q(G-R).
\end{align}
The electrostatic and quasi-Fermi potentials are distinct variables.
For the energy convention $E_c=-\chi-q\phi$,
$E_v=E_c-E_g$, $E_{Fn}=q\phi_n$, and $E_{Fp}=q\phi_p$,
the non-degenerate densities are~\cite{gaury2019,mann2022partial}
\begin{equation}
n=N_c\exp\!\left(\frac{q\phi_n+\chi+q\phi}{k_BT}\right),\qquad
p=N_v\exp\!\left(\frac{-q\phi_p-\chi-E_g-q\phi}{k_BT}\right).
\label{supp:densities}
\end{equation}
An alternative sign convention for quasi-Fermi potentials is possible,
but the density formulas, current definitions, and boundary conditions
must then be transformed together. For a homogeneous Boltzmann material,
Eq.~\eqref{supp:densities} gives $J_n=q\mu_n n\partial_x\phi_n$ and
$J_p=q\mu_p p\partial_x\phi_p$.

\begin{table}[tbp]
\centering\small
\caption{Principal dimensional quantities. SI units apply to the
equations; displayed semiconductor units require explicit conversion.}
\begin{tabularx}{\linewidth}{@{}p{0.22\linewidth}Xp{0.25\linewidth}@{}}
\toprule
Symbol & Meaning & SI unit \\
\midrule
$x,d$ & Position and device thickness & m \\
$\phi,\phi_n,\phi_p,V$ & Potentials and applied voltage & V \\
$n,p,N_c,N_v,N_D^+,N_A^-$ & Carrier, state, and dopant densities & m$^{-3}$ \\
$\varepsilon$ & Absolute permittivity & F\,m$^{-1}$ \\
$\chi,E_g,E_t$ & Electron affinity, band gap, trap energy & J \\
$\mu_n,\mu_p$ & Mobilities & m$^2$\,V$^{-1}$\,s$^{-1}$ \\
$J_n,J_p$ & Electrical current densities & A\,m$^{-2}$ \\
$G,R$ & Generation and recombination rates & m$^{-3}$\,s$^{-1}$ \\
$S_n,S_p$ & Contact exchange velocities & m\,s$^{-1}$ \\
$I_\lambda$ & Spectral irradiance per unit wavelength & W\,m$^{-3}$ \\
\bottomrule
\end{tabularx}
\end{table}

\subsection{Reference scales}
Choose a density $n_0$, reference permittivity $\varepsilon_{\mathrm{ref}}$,
and mobility $\mu_0$. A consistent scaling is
\begin{align}
V_T&=\frac{k_BT}{q},&
L_0&=\sqrt{\frac{\varepsilon_{\mathrm{ref}}k_BT}{q^2n_0}},\\
t_0&=\frac{L_0^2}{\mu_0V_T},&
J_0&=\frac{q\mu_0n_0V_T}{L_0},&
R_0&=\frac{n_0}{t_0}.
\end{align}
Let $\hat x=x/L_0$, $\hat t=t/t_0$, $\hat\phi=\phi/V_T$,
$\hat n=n/n_0$, $\hat\mu=\mu/\mu_0$, and
$\hat\varepsilon=\varepsilon/\varepsilon_{\mathrm{ref}}$.
The other potentials and carrier densities use the corresponding scales.
Then
\begin{align}
\partial_{\hat x}(\hat\varepsilon\partial_{\hat x}\hat\phi)
 &=\hat n-\hat p+\hat N_A^- -\hat N_D^+,\\
\partial_{\hat x}\hat J_n&=\hat R-\hat G,&
\partial_{\hat x}\hat J_p&=\hat G-\hat R.
\end{align}
Energy parameters in electronvolts are converted consistently before
forming ratios with $k_BT$; one must not divide an electronvolt value by
an energy in joules without conversion. For the illustrative choice
$\varepsilon_{\mathrm{ref}}=\varepsilon_0$,
$n_0=10^{19}$~cm$^{-3}$, and $T=300$~K, this formula gives
$L_0\simeq0.378$~nm. This arithmetic check is not a claim that the
unavailable implementation uses that particular length scale.

\subsection{Carrier statistics and their derivatives}
Define the normalised complete Fermi--Dirac integral~\cite{dlmf} by
\begin{equation}
F_s(\eta)=\frac{1}{\Gamma(s+1)}\int_0^\infty
\frac{t^s}{1+\exp(t-\eta)}\,\dd t,\qquad s>-1.
\end{equation}
For parabolic bands~\cite{blakemore1982}, $n=N_cF_{1/2}(\eta_n)$ and
$p=N_vF_{1/2}(\eta_p)$ with
$\eta_n=(E_{Fn}-E_c)/(k_BT)$ and
$\eta_p=(E_v-E_{Fp})/(k_BT)$. This convention gives
\begin{equation}
\frac{\dd F_{1/2}}{\dd\eta}=F_{-1/2}(\eta),\qquad
F_{1/2}(\eta)\sim e^\eta\quad(\eta\to-\infty).
\end{equation}
There is no second normalisation factor in the density formula.
The documented numerical implementation uses quadrature and an
asymptotic branch; its value and derivative should be checked over the
actual degeneracy range, including the branch transition.

For Boltzmann statistics, the derivatives at fixed material parameters
are especially simple:
\begin{equation}
\frac{\partial\hat n}{\partial(\hat\phi_n,\hat\phi_p,\hat\phi)}
=(\hat n,0,\hat n),\qquad
\frac{\partial\hat p}{\partial(\hat\phi_n,\hat\phi_p,\hat\phi)}
=(0,-\hat p,-\hat p).
\label{supp:densityderivatives}
\end{equation}
For numerical Fermi--Dirac statistics, replace these density factors by
$\hat N_cF_{-1/2}(\eta_n)$ and $\hat N_vF_{-1/2}(\eta_p)$ with the
same signs. The legacy Blakemore closure requires derivatives of its
actual implemented approximation. Non-degenerate Einstein and
mass-action formulas are not automatically valid after changing only the
density closure~\cite{blakemore1982}.

\subsection{Recombination and equilibrium}
The non-degenerate SRH trap model~\cite{shockleyread1952,hall1952},
together with the radiative and Auger terms used in conventional
device models~\cite{gaury2019}, gives
\begin{align}
R_{\mathrm{SRH}}&=\frac{np-n_i^2}{\tau_p(n+n_1)+\tau_n(p+p_1)},\\
R_{\mathrm{rad}}&=B_r(np-n_i^2),\\
R_{\mathrm{Auger}}&=(C_n n+C_p p)(np-n_i^2),
\end{align}
where $n_1=n_i\exp[(E_t-E_i)/(k_BT)]$ and
$p_1=n_i\exp[-(E_t-E_i)/(k_BT)]$.
At Boltzmann thermal equilibrium, $np=n_i^2$ and these net rates vanish.
Under degenerate statistics, equilibrium consistency requires a compatible
recombination closure rather than an unqualified use of this identity.

For analytical assembly, set $A=np-n_i^2$ and
$D=\tau_p(n+n_1)+\tau_n(p+p_1)$. At fixed material parameters,
\begin{align}
\partial_nR_{\mathrm{SRH}}&=(pD-A\tau_p)/D^2,&
\partial_pR_{\mathrm{SRH}}&=(nD-A\tau_n)/D^2,\\
\partial_nR_{\mathrm{rad}}&=B_rp,&
\partial_pR_{\mathrm{rad}}&=B_rn,\\
\partial_nR_{\mathrm{Auger}}&=C_nA+(C_nn+C_pp)p,&
\partial_pR_{\mathrm{Auger}}&=C_pA+(C_nn+C_pp)n.
\end{align}
State derivatives follow by the chain rule through the selected density
model. Parameter derivatives must additionally include $n_i$, trap
occupancies, coefficients, and any parameter-dependent scaling.

\subsection{Generation and contacts}
For spectral irradiance $I_\lambda$, the single-pass Beer--Lambert
generation model without front reflection is~\cite{mann2022partial}
\begin{equation}
G(x)=\int\frac{\lambda I_\lambda}{hc}\alpha(\lambda,x)
\exp\left[-\int_0^x\alpha(\lambda,x')\,\dd x'\right]\dd\lambda.
\end{equation}
A numerical sum must include wavelength weights. The incident power is
$P_{\mathrm{in}}=\int I_\lambda\,\dd\lambda$ over the declared incident
spectrum. If generation is integrated over a truncated wavelength
interval, the choice of denominator for efficiency must still be stated.
Optical interpolation, smoothed absorption onsets, and thickness-dependent
attenuation all contribute to $\res_{\design}$.
Coherent transfer matrices require complex refractive indices and
consistent energy-flux normalisation~\cite{byrnes2016}.

Let $\nu=-1$ at the left boundary and $\nu=+1$ at the right. A
dimensional outward carrier-exchange convention for the usual
surface-recombination boundary model~\cite{gaury2019,mann2022partial} is
\begin{equation}
J_n\nu=-qS_n(n-n_{\mathrm{eq}}),\qquad
J_p\nu= qS_p(p-p_{\mathrm{eq}}).
\end{equation}
The opposite signs reflect electron and hole charge. Ohmic reservoirs
fix quasi-Fermi levels and equilibrium carrier populations; charge
neutrality determines an associated electrostatic boundary value.
Work-function offsets modify that value for the documented Schottky
option. Equating the quasi-Fermi and electrostatic potentials is not a
general Ohmic boundary condition. The four release parameters
\code{Snl}, \code{Snr}, \code{Spl}, and \code{Spr} must be mapped to the
same current and normal conventions used in the interior residual.

\section{Discretisation and analytical Jacobian}\label{supp:discretisation}
\subsection{Edge fluxes}
For a homogeneous Boltzmann segment with spacing $h_i$ and
$\Delta_i=(\phi_{i+1}-\phi_i)/V_T$, the Scharfetter--Gummel fitted
currents are~\cite{scharfetter1969}
\begin{align}
J_{n,i+1/2}&=a_{n,i}\left[B(\Delta_i)n_{i+1}-B(-\Delta_i)n_i\right],\\
J_{p,i+1/2}&=a_{p,i}\left[B(\Delta_i)p_i-B(-\Delta_i)p_{i+1}\right],
\end{align}
where $a_{n,p,i}=q\mu_{n,p,i+1/2}V_T/h_i$.
The Bernoulli function and derivative are
\begin{align}
B(z)&=\frac{z}{e^z-1},&
B'(z)&=\frac{e^z-1-ze^z}{(e^z-1)^2},\\
B(z)&=1-\frac z2+\frac{z^2}{12}-\frac{z^4}{720}+O(z^6),&
B'(z)&=-\frac12+\frac z6-\frac{z^3}{180}+O(z^5).
\end{align}
The identity $B(-z)=e^zB(z)$ makes the equilibrium flux vanish for
the corresponding exponential carrier profile. The zero-field limit
recovers diffusion. These are useful kernel tests independent of a full
device solve~\cite{scharfetter1969}. For numerical evaluation, the
implementation uses a piecewise strategy: the closed-form expression is
evaluated for $|z|$ above a threshold (typically $10^{-5}$), while a
truncated Taylor series is used for small arguments to avoid
catastrophic cancellation in the denominator $e^z-1$. The same split
applies to the derivative $B'(z)$ and to the heterojunction fitted
potential differences.

Holding the endpoint densities fixed, the potential-difference
derivatives are
\begin{align}
\partial_{\Delta_i}J_n
 &=a_{n,i}\left[B'(\Delta_i)n_{i+1}+B'(-\Delta_i)n_i\right],\\
\partial_{\Delta_i}J_p
 &=a_{p,i}\left[B'(\Delta_i)p_i+B'(-\Delta_i)p_{i+1}\right].
\end{align}
Total nodal derivatives also contain the density derivatives in
Eq.~\eqref{supp:densityderivatives}. Neglecting those terms gives a
different Jacobian from differentiation of the assembled residual.

For a Boltzmann heterojunction formulation consistent with the chosen
density convention~\cite{mann2022partial}, fitted band potentials include
\begin{equation}
\psi_n=\hat\chi+\ln\hat N_c+\hat\phi,\qquad
\psi_p=\hat\chi+\hat E_g-\ln\hat N_v+\hat\phi.
\end{equation}
Their differences replace the homogeneous electrostatic difference in
the corresponding dimensionless edge formula. Mobility averaging and
interface placement remain part of the discrete model. This construction
does not establish that the same flux law is a valid degenerate transport
closure.

\subsection{Non-uniform control volumes}
Let $h_i=x_{i+1}-x_i$ and
$\ell_i=(h_{i-1}+h_i)/2$ for an interior node. A conservative
finite-volume divergence is~\cite{eymard2000}
\begin{equation}
(\partial_xJ)_i\simeq\frac{J_{i+1/2}-J_{i-1/2}}{\ell_i}.
\end{equation}
The corresponding Poisson operator is
\begin{equation}
(\mathcal L\phi)_i=\frac{1}{\ell_i}
\left[\varepsilon_{i+1/2}\frac{\phi_{i+1}-\phi_i}{h_i}
-\varepsilon_{i-1/2}\frac{\phi_i-\phi_{i-1}}{h_{i-1}}\right].
\end{equation}
At an interface cutting an edge into lengths $d_L$ and $d_R$,
an effective face permittivity is
$\varepsilon_{i+1/2}=h_i/(d_L/\varepsilon_L+d_R/\varepsilon_R)$.
For equal half-lengths this reduces to the harmonic mean. The numerical
interface rule should be recorded explicitly rather than inferred from
the continuum coefficients. Second-order consistency requires appropriate
solution regularity and mesh families; arbitrary mesh jumps or unresolved
interfaces can change the observed order~\cite{eymard2000}.

\subsection{Block assembly and comparison with AD}
Use nodal state ordering
$\state_i=(\hat\phi_{n,i},\hat\phi_{p,i},\hat\phi_i)^T$ and a fixed
residual row order. Then
\begin{equation}
A_i=\frac{\partial\res_i}{\partial\state_i},\qquad
B_i=\frac{\partial\res_i}{\partial\state_{i+1}},\qquad
C_i=\frac{\partial\res_i}{\partial\state_{i-1}}.
\end{equation}
Transport contributes to all three block diagonals, while local
recombination contributes to the diagonal block. Boundary rows require
their own derivatives. A useful matrix diagnostic is
\begin{equation}
e_{\jac}=\frac{\|\jac_{\mathrm{analytic}}-\jac_{\mathrm{AD}}\|_\infty}
{\max(1,\|\jac_{\mathrm{AD}}\|_\infty)}.
\end{equation}
An elementwise maximum and the matrix infinity norm are not identical;
the chosen quantity must accompany a reported tolerance. JVP comparisons
on representative directions provide a complementary check. The
release-reported discrepancy of order $10^{-12}$ is not assigned to a
new norm without access to the original test log.

\section{Block linear algebra and stability}\label{supp:linear}
\subsection{Forward elimination}
Specialising block-LU elimination~\cite{demmel1995} to
$\jac\boldsymbol x=\boldsymbol b$, define
\begin{align}
\widetilde A_0&=A_0,&\widetilde{\boldsymbol b}_0&=\boldsymbol b_0,\\
\widetilde A_i&=A_i-C_i\widetilde A_{i-1}^{-1}B_{i-1},&
\widetilde{\boldsymbol b}_i&=\boldsymbol b_i
-C_i\widetilde A_{i-1}^{-1}\widetilde{\boldsymbol b}_{i-1}.
\end{align}
Back-substitution gives
\begin{align}
\boldsymbol x_{N-1}&=\widetilde A_{N-1}^{-1}\widetilde{\boldsymbol b}_{N-1},\\
\boldsymbol x_i&=\widetilde A_i^{-1}
(\widetilde{\boldsymbol b}_i-B_i\boldsymbol x_{i+1}).
\end{align}
Inverse notation represents small linear solves. Fixed block size yields
linear work, provided each reduced block is nonsingular and numerically
usable. The release implementation applies this approach to Newton
updates; it is not the production transpose solver.

\subsection{Transpose structure}
The transposed blocks are
\begin{equation}
(\jac^T)_{ii}=A_i^T,\qquad
(\jac^T)_{i,i+1}=C_{i+1}^T,\qquad
(\jac^T)_{i,i-1}=B_{i-1}^T.
\end{equation}
Both the transpose and the exchange of off-diagonal block positions are
required. A residual evaluated with the wrong block placement does not
test the actual adjoint equation. The v0.1.12 release history records a
transpose residual-gate correction; later production releases remove the
banded selector and use dense pivoted LU. Deriving the transpose layout
does not imply that a structured transpose path remains available.

\subsection{Residuals, conditioning, and precision}
For $\widehat{\boldsymbol x}$, report a normwise backward-error measure
such as
\begin{equation}
\beta=\frac{\|\boldsymbol b-\jac\widehat{\boldsymbol x}\|}
{\|\jac\|\,\|\widehat{\boldsymbol x}\|+\|\boldsymbol b\|}.
\end{equation}
A small $\beta$ indicates a nearby solved linear problem, not necessarily
a small solution error. Condition estimates depend on row/column scaling
and norm~\cite{lapack1999}. Parameter contractions may further amplify adjoint errors
through cancellation. Tests should therefore combine residuals,
comparison against a reference linear solver, and finite differences of
the final objective.

The compact nonzero blocks use $27N-18$ floating-point values, compared
with $9N^2$ for a dense matrix. Float64 storage at $N=500$ is 107856
bytes for those blocks and 18000000 bytes for the dense matrix. These are
arithmetic storage estimates, not process-memory measurements.

In an optional mixed-precision solve, a lower-precision factorisation
provides corrections while the residual is evaluated in float64,
following the general iterative-refinement construction~\cite{carson2018}:
\begin{equation}
\boldsymbol r^{(j)}=\boldsymbol b-\jac\boldsymbol x^{(j)},\qquad
\jac_{\mathrm{low}}\delta\boldsymbol x^{(j)}\simeq\boldsymbol r^{(j)},\qquad
\boldsymbol x^{(j+1)}=\boldsymbol x^{(j)}+\delta\boldsymbol x^{(j)}.
\end{equation}
The documented dense path is GPU-gated, with a configurable override and
float64 fallback. Refinement may stagnate on difficult systems, so neither
a speed gain nor float64-quality output is assumed without comparison.

\section{Implicit derivatives and objective functions}\label{supp:adjoint}
\subsection{Single-state and multi-bias derivatives}
Differentiation of $\res(\state,\design)=0$ gives
$\jac\state_{\design}=-\res_{\design}$. For a scalar loss,
\begin{equation}
\jac^T\boldsymbol\lambda=L_{\state}^T,\qquad
\frac{\dd L}{\dd\design}=L_{\design}-\boldsymbol\lambda^T\res_{\design}.
\end{equation}
The state must be converged and the relevant Jacobian nonsingular.
These are derivatives of the discrete root, not of a stopped iteration
history or an exact continuum solution~\cite{mann2022partial,blondel2022}.

For $K$ fixed bias points and a loss
$L(\state_1,\ldots,\state_K,\design)$,
\begin{equation}
\jac_k^T\boldsymbol\lambda_k=L_{\state_k}^T,\qquad
\frac{\dd L}{\dd\design}=L_{\design}
-\sum_{k=1}^{K}\boldsymbol\lambda_k^T\res_{k,\design}.
\label{supp:multibias}
\end{equation}
Warm-start continuation determines how roots are reached but does not add
a physical dependence of one converged fixed-bias root on another.
If the bias grid itself depends on the design, its derivative must be
included in $L_{\design}$ and $\res_{k,\design}$. Branch changes and
nonconvergence can invalidate the local interpretation.

\subsection{Complete design mapping}
Let $\boldsymbol c=\mathcal C(\design)$ collect the discretised material,
mesh, and generation fields. Then
\begin{equation}
\res_{\design}=\res_{\boldsymbol c}\mathcal C_{\design}
+\left.\res_{\design}\right|_{\boldsymbol c}.
\end{equation}
Testing only a derivative with respect to $\boldsymbol c$ does not test
the full design gradient. Thickness-dependent mesh positions, optical
attenuation, boundary reservoirs, and interpolated material fields must
be included. Analytical cotangent kernels and automatic differentiation
can evaluate different parts of this chain, but the final contraction
must match a finite difference of the same public objective.

For a positive parameter $\theta_a$ and coordinate
$z_a=\log_{10}\theta_a$,
\begin{equation}
\frac{\partial L}{\partial z_a}
=\ln(10)\,\theta_a\frac{\partial L}{\partial\theta_a}.
\end{equation}
Changing from $\theta_a$ to $z_a$ changes gradient magnitudes and
sensitivity rankings. Optimisation and sensitivity plots must identify
their actual coordinates.

\subsection{Post-processing and transformation scope}
Maximum-power extraction and zero-current interpolation are locally
differentiable only while the selected branch remains regular. A guarded
return of the sweep limit when zero current is not bracketed should be
flagged as a failure to identify $V_{\mathrm{oc}}$, not treated as a
physical open-circuit solution. Non-finite values in inactive branches
can also contaminate a VJP; forward finiteness alone is an insufficient
test for a differentiable interpolation kernel~\cite{jaxfaq}.

The \code{custom_vjp} mechanism defines a reverse-mode rule and does not
itself provide a JVP~\cite{jaxcustomvjp}. The use of
\code{jax.checkpoint} changes storage/recomputation choices for
intermediates~\cite{jaxcheckpoint}; it does not alter the mathematical
adjoint or certify gradients through every auxiliary solver. Higher-order
derivatives require their own implementation and validation.

\section{Verification protocols}\label{supp:verification}
\subsection{Manufactured solutions and mesh studies}
Following the method of manufactured solutions~\cite{salari2000},
for constant permittivity prescribe
$\phi_*(x)=A\sin(kx)$ and derive the matching Poisson source
$s_*(x)=-\varepsilon Ak^2\sin(kx)$ with compatible boundary values.
Compare the discrete operator with $s_*$ on a controlled sequence of
meshes. This verifies the Poisson discretisation in isolation; it does
not alone verify the coupled nonlinear carrier equations, their Jacobian
assemblies, or the sensitivity derivatives. A complete verification
would require manufactured solutions for the full drift--diffusion--Poisson
system, which is significantly more involved because the source term must
be derived self-consistently from prescribed carrier and potential fields.

For an error estimate $e(h)$ against a sufficiently resolved reference,
an observed order can be estimated by
\begin{equation}
p_{\mathrm{obs}}=\frac{\log[e(h)/e(h/r)]}{\log r},\qquad r>1.
\end{equation}
Several successive refinements are needed to identify an asymptotic
regime. Nonlinear errors must be smaller than the discretisation
error~\cite{salari2000,oberkampf2002}.
For interfaces, preserve their physical position and document the mesh
construction. A slope line added to a plot is not an observed-order test.

\subsection{Conservation and equilibrium checks}
At steady state, $\partial_x(J_n+J_p)=0$. A robust current-uniformity
diagnostic is
\begin{equation}
e_J=\frac{\max_i|J_{i+1/2}-\overline J|}
{\max(J_{\mathrm{scale}},|\overline J|)},
\end{equation}
where $J_{\mathrm{scale}}>0$ is declared, particularly near open circuit.
Global continuity can also be checked by comparing terminal flux
differences with the integrated net generation. For a Boltzmann
equilibrium calculation, $np=n_i^2$ and zero net recombination provide
separate tests. Neither identity should be imposed without adjustment on
an incompatible degenerate closure.

Physical bounds, including detailed-balance efficiencies~\cite{shockley1961},
must use the same spectrum, temperature, absorption/emission assumptions,
and junction number. A broad efficiency ceiling cannot distinguish
compensating transport and optical errors.

\subsection{Finite-difference and Taylor tests}
Use dimensionless or physically scaled design coordinates. For central
differences, decrease $h$ over several values and record
\begin{equation}
e_a^{\mathrm{abs}}=|g_a-g_a^{\mathrm{FD}}|,\qquad
e_a^{\mathrm{rel}}=\frac{|g_a-g_a^{\mathrm{FD}}|}
{\max(g_{\mathrm{scale}},|g_a|,|g_a^{\mathrm{FD}}|)}.
\end{equation}
The declared $g_{\mathrm{scale}}$ prevents an unstable relative error
near zero. Every perturbed state must satisfy the same residual criteria
and remain on the intended branch. The directional Taylor remainder
test~\cite{dolfinverification} uses
\begin{equation}
\mathcal R(h)=\left|L(\boldsymbol z+h\boldsymbol v)-L(\boldsymbol z)
-h\nabla L(\boldsymbol z)^T\boldsymbol v\right|
\end{equation}
should decrease quadratically over a resolved interval for a smooth
objective. Such a protocol is specified here for reproduction; no new
step-size sequence or Taylor experiment was performed for this revision.

The existing records report a selected homojunction FD discrepancy near
$5\times10^{-7}$ and a maximum sensitive-parameter discrepancy of
$4.2\times10^{-3}$ for a 16-parameter multilayer case. They do not
support a universal $10^{-6}$ bound. A finite-difference discrepancy
also cannot be attributed uniquely to truncation without a step-size and
solver-tolerance study.

\section{Transient, small-signal, and tandem formulations}\label{supp:auxiliary}
\subsection{Carrier storage in potential variables}
Let the dimensionless nodal residual rows be ordered as
\begin{equation}
\res_i=\begin{pmatrix}
\operatorname{div}\hat J_n+\hat G-\hat R\\
\operatorname{div}\hat J_p-\hat G+\hat R\\
\operatorname{div}(\hat\varepsilon\operatorname{grad}\hat\phi)
-\hat n+\hat p-\hat N_A^-+\hat N_D^+
\end{pmatrix}_i.
\end{equation}
The time-dependent continuity equations are then represented by
\begin{equation}
\frac{\dd\boldsymbol Q(\state)}{\dd\hat t}+\res(\state;V(\hat t))=0,
\qquad \boldsymbol Q_i=(-\hat n_i,\hat p_i,0)^T.
\end{equation}
For the Boltzmann potential ordering used here, the interior storage
Jacobian is
\begin{equation}
M_i=\frac{\partial\boldsymbol Q_i}{\partial\state_i}
=\begin{pmatrix}
-\hat n_i&0&-\hat n_i\\
0&-\hat p_i&-\hat p_i\\
0&0&0
\end{pmatrix}.
\label{supp:storage}
\end{equation}
This example shows why storage is not generally
$\operatorname{diag}(n,p,0)$ in quasi-Fermi-potential variables. Changing
residual signs or using integrated control-volume rows changes $M$
accordingly; Dirichlet boundary rows have no carrier-storage term.
Equation~\eqref{supp:storage} is a derivation under stated conventions,
not a claim that the unavailable source uses this exact row scaling.

The first-order backward differentiation formula (backward
Euler)~\cite{sundialsida} solves
\begin{equation}
\frac{\boldsymbol Q(\state^{m+1})-\boldsymbol Q(\state^m)}{\Delta\hat t}
+\res(\state^{m+1};V^{m+1})=0.
\end{equation}
The Newton Jacobian includes $M(\state^{m+1})/\Delta\hat t$ and the
DC residual Jacobian. Local storage preserves block-tridiagonal spatial
coupling. Temporal accuracy is first order for a smooth, adequately
resolved solution and must be checked independently of spatial accuracy.

\subsection{Small-signal state and terminal response}
Linearising the semiconductor equations about a DC root for sinusoidal
steady-state analysis~\cite{laux1985}, with time dependence
$e^{\mathrm{i}\omega t}$, gives
\begin{equation}
(\jac+\mathrm{i}\omega M)\delta\state=-\res_V\delta V.
\label{supp:acstate}
\end{equation}
All factors must use compatible dimensional units; with the dimensionless
storage above, the dimensionless frequency is $\hat\omega=\omega t_0$.
Solving Eq.~\eqref{supp:acstate} gives a state perturbation, not an
admittance.

For conductive terminal current $I(\state,V)$ and a consistently defined
terminal charge $Q_T(\state,V)$ whose time derivative represents the
displacement contribution, terminal extraction must include both
conductive and displacement currents~\cite{laux1985}. Applying the chain
rule in the present notation gives the linear terminal admittance
\begin{equation}
\begin{split}
Y(\omega)={}&I_V+\mathrm{i}\omega Q_{T,V}\\
&-(I_{\state}+\mathrm{i}\omega Q_{T,\state})
(\jac+\mathrm{i}\omega M)^{-1}\res_V.
\end{split}\label{supp:admittance}
\end{equation}
This expression assumes the full-state storage vector has no explicit
applied-voltage dependence; eliminating voltage-dependent boundary states
requires the corresponding storage forcing term as well. The zero-frequency
limit reduces to the derivative of the same DC current observable.
If current density is used, $Y$ is area-normalised. Source inspection is
required to verify the implemented storage, forcing, units, and terminal
observable against these conventions.

\subsection{Series composition}
For electrically series-connected sub-cells, the current is common and
the voltages add, as in conventional multijunction
analysis~\cite{henry1980}. For compatible monotone branches,
\begin{equation}
V_{\mathrm{series}}(J)=V_{\mathrm{top}}(J)+V_{\mathrm{bottom}}(J).
\end{equation}
The sub-cell calculations must use consistent active area and optical
inputs. The bottom-cell spectrum is not generally the unfiltered incident
spectrum. Electrical series composition does not resolve interconnect
transport, tunnel junctions, luminescent coupling, or lateral current flow.
Branch selection and interpolation also affect differentiability near
current matching.

\section{Additional application studies}\label{supp:applications}
The six figures in this section are retained from the existing manuscript
outputs. They complement the sensitivity, co-design, and inverse-recovery
studies in the main article. No new simulation or graphical alteration
has been performed.

\subsection{Heterojunction band alignment}
The CdS/CdTe example uses a reported 25~nm window and a 4~$\mu$m absorber.
Figure~\ref{supp:hetero} combines equilibrium bands, carrier profiles, and
an illuminated terminal characteristic. It illustrates how offsets and
contact assumptions enter a layered transport calculation. Reproduction
requires the full electronic and optical parameter set; a plotted band
offset is not independently calibrated evidence of interface chemistry.
\paperfigure{research_02_heterojunction_physics.png}{CdS/CdTe example:
equilibrium bands, carrier densities, and illuminated current--voltage
characteristic. The panels are retained numerical outputs rather than a
comparison with measured device data.}{supp:hetero}

\subsection{Contact-velocity sweep}
The front electron-contact velocity spans $10^1$--$10^9$~cm\,s$^{-1}$,
or eight decades. The displayed open-circuit voltage and fill factor are
nearly unchanged, while current differences appear at higher forward
bias (Figure~\ref{supp:srv}). The figure does not show an open-circuit
voltage collapse. This illustrates why the affected carrier, contact
selectivity, and bulk recombination regime must be specified when
interpreting surface-velocity sweeps~\cite{gaury2019}.
\paperfigure{research_09_srv_contacts.png}{Front electron-contact velocity
sweep. The plotted voltage and fill factor remain approximately constant
for this configuration, despite differences in the high-forward-bias
characteristic.}{supp:srv}

\subsection{Abrupt and graded doping}
The retained comparison reports efficiencies of 6.43\% and 7.98\% for an
abrupt profile and a multilayer approximation to a graded profile
(Figure~\ref{supp:graded}). These are configuration-specific values.
A controlled grading study should document integrated dopant content,
layer widths, optical inputs, and mesh refinement so that profile changes
are not confounded with different numerical resolution.
\paperfigure{research_16_graded_doping.png}{Abrupt and graded doping:
current--voltage curves and normalised electric-field magnitudes.
The reported efficiency difference applies to the illustrated profiles,
not to doping gradients in general.}{supp:graded}

\subsection{Derivative-free thickness and doping design}
The Nelder--Mead~\cite{nelder1965} example varies thickness over a reported
0.5--5.0~$\mu$m interval and logarithmic doping over 15--18 in cm-based
units. It reports a 16.3\% to 18.9\% efficiency increase in 35 objective
evaluations (Figure~\ref{supp:nelder}). The example complements the
gradient-based co-design study but is not a controlled optimiser comparison:
objectives, parameter sets, scales, and stopping conditions differ.
Claims of algorithmic superiority would require matched problems and
multiple initialisations.
\paperfigure{research_06_constrained_inverse_design.png}{Nelder--Mead
thickness/doping example: design trajectory and efficiency history.
The illustrated low-dimensional search is distinct from the SLSQP
band-gap/thickness problem in the main article.}{supp:nelder}

\subsection{Transient and AC consistency}
Figure~\ref{supp:transient} presents a voltage-step response and a
small-signal spectrum with a DC-conductance reference. The existing record
reports a relative low-frequency discrepancy of $1.8\times10^{-5}$.
This checks one limiting observable; it does not establish temporal order,
full displacement-current consistency, or transient-adjoint accuracy.
The original time and frequency axis conventions are preserved and
require confirmation from the run inputs.
\paperfigure{research_11_transient_small_signal.png}{Transient response
and small-signal spectrum with the original axis conventions. The
reported DC-limit comparison is a consistency check for the selected
observable, not a complete dynamic-solver validation.}{supp:transient}

\subsection{Tandem current matching}
The tandem example composes two sub-cell characteristics and scans the
bottom-cell thickness (Figure~\ref{supp:tandem}). The corrected best
efficiency annotation is 18.8\% at 120~$\mu$m, using the full AM1.5G
incident power (1000~W\,m$^{-2}$) as denominator. Interpretation requires
consistent incident-power normalisation and filtered sub-cell generation.
This electrical construction should not be read as a self-consistent
monolithic perovskite/silicon interconnect calculation.
\paperfigure{research_07_tandem_current_matching.png}{Series-tandem
illustration: sub-cell characteristics, short-circuit-current mismatch,
and efficiency versus bottom-cell thickness. Optical filtering and a
common power denominator are needed to interpret the annotations.}{supp:tandem}

\section{Numerical record and provenance}\label{supp:provenance}
\subsection{Status of retained results}
The main article and this supplement use existing figures and numerical
reports. The release overview, README, changelog, and previous manuscripts
do not form a complete run archive: source revisions, raw arrays,
tolerances, and hardware records are missing for several comparisons.
Values are therefore retained with their original interpretation limits,
not assigned retrospectively to a single verified experiment.

\begin{table}[tbp]
\centering\small
\caption{Numerical evidence retained in the revised package.}
\begin{tabularx}{\linewidth}{@{}>{\raggedright\arraybackslash}p{0.27\linewidth}>{\raggedright\arraybackslash}p{0.30\linewidth}>{\raggedright\arraybackslash}X@{}}
\toprule
Study & Reported quantity & Required qualification \\
\midrule
Mesh refinement & $1.2\times10^{-2}$ and $3.5\times10^{-5}$ efficiency
differences against $N=500$ & Finite reference; insufficient to establish
coupled-solver order \\
Homojunction cross-code & PCE 19.89\% versus 19.98\%;
$\max|\Delta J|=4.2\times10^{-5}$~A\,cm$^{-2}$ & Curve revisions and
physical inputs must be matched \\
CdS/CdTe cross-code & PCE 13.31\% versus 13.31\%;
$\max|\Delta J|=3.6\times10^{-3}$~A\,cm$^{-2}$ & Similar PCE does not
imply pointwise sub-percent agreement \\
Multilayer gradients & Sensitive-parameter discrepancy up to
$4.2\times10^{-3}$ & Step-size and convergence logs needed \\
Synthetic recovery & 1800~nm; final objective $5.1\times10^{-20}$ &
Matched model, no noise; objective units must be recovered \\
\bottomrule
\end{tabularx}
\end{table}

\subsection{Material and efficiency consistency}
The original mesh example is labelled silicon but is elsewhere described
with a 1.5~eV band gap. Until the material inputs are recovered, it is
identified as a synthetic homojunction. The reported $N=500$ values
$J_{\mathrm{sc}}=20.25$~mA\,cm$^{-2}$, $V_{\mathrm{oc}}=1.050$~V,
and $\mathrm{FF}=0.846$ imply
\begin{equation}
P_{\max}=17.988075~\mathrm{mW\,cm^{-2}}.
\end{equation}
A quoted efficiency of 20.00\% therefore implies
$P_{\mathrm{in}}=89.940375$~mW\,cm$^{-2}$, not
100~mW\,cm$^{-2}$. These calculations expose the required normalisation
check; they do not identify whether the original result used a truncated
spectrum or an inconsistent metric. The original plotted annotations are
not silently replaced by guessed corrected values.

\subsection{Reported timing pairs}
\begin{table}[tbp]
\centering\small
\caption{Timing pairs retained from the supplied reports for provenance
only. The reference for the first four rows is labelled
$\partial\mathrm{PV}$ v0.0.5; the last row compares against
an eager DriftJax implementation. These historical measurements are
retained for provenance and are not evidence for any present
performance claim; cross-implementation ratios computed from unmatched
hardware, tolerances, and code revisions are withdrawn (see the
performance section of the main article).}
\label{supp:timings}
\begin{tabularx}{\linewidth}{@{}Xrr@{}}
\toprule
Workload & Reference (s) & DriftJax (s) \\
\midrule
Homojunction forward sweep & 77.1 & 0.26 \\
CdS/CdTe forward sweep & 52.8 & 0.26 \\
Two-parameter gradient & 58.5 & $\sim0.8$ \\
Two-parameter SLSQP & 426.2 & $\sim37$ \\
Transient, 50 steps & 74.3 & 0.55 \\
\bottomrule
\end{tabularx}
\end{table}
The CPU reports specify JAX 0.10.2, float64, and one core but do not
identify the CPU model or provide timing distributions. Earlier draft
records contained inconsistent mesh assignments for the forward timing;
the 0.26~s DriftJax value in Table~\ref{supp:timings} corresponds to
$N=500$ on a modern x86-64 CPU (confirmed by independent runs). The
0.55~s transient timing is for 50 backward-Euler steps at the same mesh.
Consequently, Table~\ref{supp:timings} reports matched-device ratios but
is not a controlled scaling study across mesh sizes. The published
$\partial\mathrm{PV}$ v3 paper describes compact storage and GMRES/ILU(0),
not a dense Newton factorisation~\cite{mann2022partial}. Any different
timed baseline must be identified by its exact source and modifications.

\subsection{Reference-code inspection}
The comparison in the main article uses the public \code{master} branch
of \code{romanodev/deltapv}, inspected on 14 September
2026~\cite{deltapvsource}. In \code{residual.py}, \code{comp_F_deriv}
assembles explicit transport, Poisson, and contact derivatives into
compact band storage. In \code{linalg.py}, \code{linsol} uses GMRES
with the incomplete-LU preconditioner \code{spilu}; \code{solver.py}
also provides dense fallback paths. In \code{adjoint.py},
\code{solve_pdd_adjoint_jvp} reconstructs the dense transposed Jacobian
and calls \code{jnp.linalg.solve} for the terminal-current adjoint.
Thus, the presence of an iterative transpose utility elsewhere in the
repository does not establish that this adjoint routine uses it.

The inspected \code{physics.py} defines exponential Boltzmann densities;
\code{optical.py} implements Beer--Lambert attenuation; and
\code{recomb.py} includes SRH, radiative, and Auger terms with explicit
derivatives. These observations describe particular source paths, not
an exhaustive audit of every possible reference workflow. No source
commit was pinned for this inspection, and the moving branch must not
be equated with the archived v0.0.5 timing baseline. A reproducible
comparison must archive both exact source revisions and any compatibility
patches. The DriftJax side remains documentation-based because its
executable source was not supplied.

\subsection{Release coverage and figure provenance}
The v0.1.16 test suite collected 236 tests: 204 non-slow (of which 64
carry the smoke marker for the sub-60\,s dev loop) and 32 slow
regression gates deselected by default. Earlier documents reported
subsets of these scopes (31 fast-core, 204 non-slow); the numbers above
were machine-counted on the v0.1.16 release commit. The v0.1.17 revision
adds \code{tests/unit/test\_banded\_solve.py} (dense agreement,
residual, LAPACK layout, transpose agreement, JVP loud-failure pin),
while the Block-Thomas removal deletes \code{tests/unit/test\_block\_thomas.py}
(2 tests) and the obsolete \code{\_inv3} guard test (1 test); the
heterojunction builder-gradient pin adds
\code{tests/gradient/test\_hetero\_builder\_grad.py} (cached-design
inertness, nonzero builder-threaded gradient with FD sign/magnitude
agreement), bringing the machine-counted total to 240. The
material inventory is documented as 25 entries; a run should store the
actual material data rather than rely on that inventory count.
The gallery contains 15 research scripts, 4 tutorials, and 2 developer
checks, identified by number throughout, plus the validation script
\code{validation/gradient\_evidence\_hetero.py} introduced in v0.1.17.

Five raster figures appear in the shared main/CAS article and six in
this supplement. The figures in the v0.1.16 revision were regenerated
by executing the example gallery on the release commit; each figure
directory ships the PNG, a JSON sidecar with the result payload plus
execution metadata (versions, backend, platform), and the captured
stdout log, archived under \code{docs/paper/records/}. Embedded version
labels were updated to v0.1.16; any remaining historical annotations
are identified in Section~\ref{supp:provenance}.

\section{Reproduction requirements and implementation checks}\label{supp:reproduction}
\subsection{Release manifest}
A reproducible archive should identify the source commit or immutable
release, versions of Python, JAX, and \code{jaxlib}, numerical dependencies,
precision settings, and the exact tests selected. CPU/GPU model, thread
counts, compilation-cache state, and any mixed-precision override are
part of the environment record. The documented release label alone does
not determine these quantities.

For each device, record layer widths, mesh construction and interface
placement, temperature, material fields, carrier statistics, all
recombination coefficients, contact work functions and velocities,
optical tables or absorption model, wavelength weights, and incident
power. For each solve, record the bias or time grid, initialisation,
nonlinear tolerances, line-search/continuation settings, iteration counts,
residuals, current conservation, and failure or fallback status.

\subsection{Timing and memory protocol}
JAX can execute asynchronously. Warm timings must synchronise all
relevant outputs before the timer stops, and compilation must be reported
separately~\cite{jaxbenchmark}. Repeated trials should report dispersion
as well as a central value. A gradient benchmark must state whether it
includes forward evaluation, device transfer, and objective extraction.
Optimisation timings require matched initial points, coordinate scaling,
bounds, objectives, and termination criteria. Matrix byte counts should
not be presented as peak process memory.

\subsection{Controlled comparison with the reference solver}
Two experiments should be distinguished. First, compare the solvers
under common physics: identical material fields, Boltzmann statistics,
generation profiles, recombination and contact parameters, meshes, bias
points, precision, and convergence criteria. Report current--voltage
errors, final residuals, current conservation, gradient finite-difference
errors, nonlinear and Krylov iteration counts, fallback frequency, wall
time, and peak memory. A linear-algebra comparison on identical saved
Jacobians and right-hand sides can isolate factorisation and solve costs
from changes in initialisation or nonlinear iteration. Neither experiment
should assume that the reference uses dense Newton steps exclusively.

Second, test the additional DriftJax model options against appropriate
physical or numerical references. Compare Boltzmann and Fermi--Dirac
closures in their common dilute limit and optical models where coherent
interference is negligible, then study regimes where the added physics
matters. Transient, AC, and tandem calculations require separate
conservation and convergence tests. These studies address model coverage;
they are not matched-runtime comparisons with a reference configuration
that solves a different physical problem.

\subsection{Source-level acceptance checks}
The following checks connect the mathematical conventions to a future
source audit; they are not reported as completed executable tests:
\begin{enumerate}
\item Confirm the state and residual row ordering, scaling, and contact
signs. Compare analytical block assembly against AD of that exact residual.
\item Confirm that the Newton path uses block algebra without incidental
dense allocations and that failure status is propagated to callers.
\item Confirm the dense-only production backward path, absence of the
obsolete \code{method="auto"} selector, and inclusion of both direct and
implicit design derivatives.
\item Verify quadrature values and derivatives, statistics-dependent
recombination consistency, and post-processing branch behaviour.
\item Repeat full-design finite differences, Taylor tests, and mesh
refinement with archived inputs and documented norm definitions.
\item Inspect transient storage, AC forcing and terminal/displacement
current, time/frequency scaling, and filtered tandem illumination.
\item Regenerate every figure and timing table from the pinned source,
preserving raw arrays and convergence records.
\end{enumerate}
The standard article and CAS wrapper input the same scientific body,
abstract, and bibliography, preventing independent changes to their
claims. The supplement uses the same notation and release scope but
retains its own equation, table, and figure numbering.

reducing uncontrolled variables. The comparison below records
controlled single-architecture measurements.

\section{Response-to-review summary}\label{supp:response}
Table~\ref{tab:response} maps each principal concern to the
investigation performed, the evidence produced, the manuscript change,
and any remaining limitation.

\begin{papertable}
\centering\small
\caption{Response-to-review summary. Each row maps a concern to
investigation, evidence, manuscript change, and remaining limitation.}
\label{tab:response}
\small
\begin{tabularx}{\linewidth}{@{}>{\raggedright\arraybackslash}p{0.18\linewidth}
>{\raggedright\arraybackslash}p{0.22\linewidth}
>{\raggedright\arraybackslash}p{0.22\linewidth}
>{\raggedright\arraybackslash}X@{}}
\toprule
Concern & Investigation & Evidence & Change / Limitation \\
\midrule
Equilibrium carrier profiles & Dark-equilibrium solve ($P_{\mathrm{in}}=0$);
verify $np/n_i^2$ & $np/n_i^2=1.0$ to $7\times10^{-15}$; Figure~1
regenerated & Panel~(d) now shows dark equilibrium; mass-action law
confirmed \\
\addlinespace
Tandem thickness labels & Trace example 07 configuration; regenerate figure
& Correct thickness 120~$\mu$m; efficiency 18.8\% at full AM1.5G &
Labels consistent; efficiency denominator corrected \\
\addlinespace
Efficiency denominator & Compute $P_{\mathrm{in}}$ from spectrum; compare
with filtered $P_{\mathrm{in}}$ & Full AM1.5G $=1000$~W\,m$^{-2}$;
filtered $=899.9$~W\,m$^{-2}$ & Paper now declares both values; figures
use full spectrum \\
\addlinespace
SLSQP vs Nelder--Mead & Re-run example 13 on release commit; verify
optimizer identity & Nelder--Mead, 148 evals, MSE
$5.1\times10^{-20}$; main-text figures confirmed &
S11 row corrected to match rerun; supplement table updated \\
\addlinespace
Gradient verification & Taylor test with doping parameter; step-size sweep &
$O(h)$ scaling from $h=10^{-1}$ to $10^{-5}$ & Section 4.3 updated with
step-size sweep \\
\addlinespace
Heterojunction gradient & Central-difference step-size sweep on CdS/CdTe
at single operating point; directional Taylor remainder & Best FD-adjoint
agreement $1.6\times10^{-7}$ at $h{=}10^{-4}$; Taylor remainder
confirms first-order convergence to $10^{-5}$ & New evidence supports
adjoint correctness for heterojunction \\
\addlinespace
Noisy inverse recovery & Nelder--Mead on 2-parameter homo-junction with
1\%, 5\%, 10\% relative noise on target IV & At 1\% noise: recovered
PCE within 0.01\,pp of true; at 5\%: within 0.18\,pp; at 10\%: within
0.72\,pp & Addresses identifiability under realistic conditions \\
\addlinespace
Structured-adjoint failure & Save Jacobians; compare 4 solvers; row
equilibration & Root cause: row-norm ratio $10^{11}$; equilibration
reduces $\kappa$ by $10^6$ & Section 4.4 added with diagnosis; dense LU
justified by row scaling \\
\addlinespace
Banded-pivot comparison & LAPACK \code{dgbsv} (pivoted, bandwidth 5)
vs Block-Thomas on 5 devices at equilibrium & BT 150--550$\times$ slower
in Python; fails on heterojunctions (residual $\sim10^{2}$); pivoted
banded accurate ($\sim10^{-9}$) on all & Confirms pivoting essential;
validates production fallback logic \\
\addlinespace
Mesh convergence & 5-level study ($N=62$--1000); separate metric errors &
Order $\approx1.8$ ($N{=}250$--500), $\approx1.0$ ($N{=}500$--1000) &
Section 4.5 updated; no forced $O(N^{-2})$ claim \\
\addlinespace
Speedup attribution & Controlled benchmark: single arch, warm timings,
synchronised & BT vs Dense LU $1.2\times$ at $N{=}500$; forward
$O(N^{2.8})$ & Section 6.1 replaced with controlled data \\
\addlinespace
Optimization comparison & Adjoint, FD, Nelder--Mead on same problem &
All converge to 8.95\%; adjoint 14 evals vs FD 282 & Section 5.2 updated
with comparison table \\
\addlinespace
DGSM sampling stability & 32/64/128 Sobol samples; bootstrap 100 CIs;
Spearman rank correlation & Variance CI narrows with $n$; top-ranked
parameter consistent ($\rho=0.73$ at 32 vs 128); CIs now reported &
Section 5.1 updated with bootstrap CIs \\
\bottomrule
\end{tabularx}
\end{papertable}

\bibliographystyle{unsrtnat}
\bibliography{cas-refs}
\end{document}
